% Source-supported replacement paragraph for infinite_status.tex.
% This file is an audit artifact, not a mutation of the mathematical source.
For the quantitative quotient extraction, retain the fixed source
trajectory with its fixed low-order tolerance chosen sufficiently small
before any frequencies, as in source pp.~12 and 69. Put
\[
 \mathcal G_q=-Je_q\cdot G_{<q},\qquad
 d_q=r_q\sin(\phi_q-\phi_{q-1}),\qquad
 a_q=\frac{\mathcal G_q}{K_qr_q},\qquad
 b_q=K_q\zeta_{q,1}.
\]
These are the physical-coordinate versions of (3.55)--(3.56),
p.~26. On growth, the preceding phase is vertical by the
holding equation, so $\zeta_{q,1}=d_q$ and $b_q=K_qd_q$.
The bounds at the end of Subsection~3.7.4 and the beginning of
Subsection~3.7.5, p.~29, together with (3.63), give, for one
fixed $0<\delta\leq1/4$,
\[
 r_q\geq1-\delta,\quad
 \frac{d_q}{\sin s_q}\geq1-\delta,\quad
 \frac{\mathcal G_q}{\sigma_{q-1}^2\sin s_q}\leq1+\delta,
 \quad
 |(\log u_*)'|\leq\delta\gamma,
 \qquad u_*:=\sqrt{b_q/a_q},\quad\gamma:=\sqrt{a_qb_q}.
\]
All these errors are fixed low-order errors: the derivative
estimate is $C\epsilon_{\rm sch}\Gamma_q/L_q$, while
$\gamma=(1+O(\epsilon_{\rm sch}))\Gamma_q$.
Their finite collection is absorbed by the same fixed tolerance
choice; this choice precedes the frequency schedule. Source
Lemma~3.3, p.~14, applies with its local tolerance renamed
$\delta$ to avoid confusing it with the physical time-scaling
parameter $\nu$.

Before any possible zero of $\Theta_q$, set
$u=\Omega_q/\Theta_q$ and $z=u/u_*$. The source entry datum
gives $z(t_q^{\rm in})=1$, and differentiation gives
\[
 u'=b_q-a_qu^2,\qquad
 z'=\gamma(1-z^2)-z(\log u_*)'.
\]
At the upper boundary $z=1+\delta$, this derivative is at most
$\gamma[1-(1+\delta)^2+\delta(1+\delta)]
=-\delta\gamma<0$.
At the lower boundary $z=1-\delta$, it is at least
$\gamma[1-(1-\delta)^2-\delta(1-\delta)]
=\delta\gamma>0$.
Thus $1-\delta\leq z\leq1+\delta$ by the first-exit argument.
The coefficient functions are continuous on the compact growth
interval and $a_q,b_q$ are positive, so $u_*$ is bounded there.
If $\Theta_q$ first vanished, the bounded quotient would force
$\Omega_q$ to vanish at the same time. Backward uniqueness of
the continuous-coefficient linear amplitude system would then
contradict its nonzero entry datum. The quotient is therefore
defined through the crossing. This is exactly the continuation
step used in the source growth proof on p.~30.

The coefficient identity now gives
\[
 u_*=K_q\sqrt{\frac{d_qr_q}{\mathcal G_q}}
 \geq\frac{K_q}{\sigma_{q-1}}
 \frac{1-\delta}{\sqrt{1+\delta}},\qquad
 \left|\frac{\Omega_q}{\Theta_q}\right|
 \geq\frac{K_q}{\sigma_{q-1}}
 \frac{(1-\delta)^2}{\sqrt{1+\delta}}
 >\frac{K_q}{2\sigma_{q-1}}.
\]
Indeed the last scalar factor decreases on $[0,1/4]$ and its
value at $1/4$ is $9/(8\sqrt5)>1/2$, since $81>80$.
This proves the bound for $q\geq2$. For $q=1$, source p.~27
gives the exact growing eigenline; its physical form is
$|\Omega_1/\Theta_1|=K_1/\nu=K_1/\sigma_0$ and $r_1=1$,
which proves the same bound.
